% Source-aligned English for printed p.207.
\clearpage
\SourcePage{207}

\LayerTarget{tohoku:remark:smith-fixed-point-application:p1}
\noindent in [2] is entirely unnecessary. As another easy application of the
preceding corollary, if a finite group $G$ of prime order $p$ acts on a
finite-dimensional space $X$ which is acyclic mod. $p$, then the fixed-point
set $X_0$ and the quotient space $X/G$ are acyclic mod. $p$. It follows that
if $X$ is acyclic for a coefficient field $k$ (of arbitrary characteristic,
the case of characteristic $\ne p$ being trivial), then so is $X/G$; and,
since $k$ is arbitrary, this statement remains valid with $k$ replaced by
$\mathbf Z$. An immediate induction shows that these results concerning $X/G$
remain valid if one assumes only that $G$ is finite and solvable. Finally, the
calculations of this section are also useful in studying Steenrod powers in
sheaves and the cohomology of symmetric powers of a space. Indeed, analogous
calculations in an unpublished manuscript of R. Godement on Steenrod powers
inspired the present section.

\LayerTarget{tohoku:section:5.5}
\bigskip
\textbf{5.5. Calculation of $H^n(X;G,A)$ by coverings.} We shall
need the following auxiliary results.

\LayerTarget{tohoku:lemma:5.5.1}
\textsc{Lemma 5.5.1.} \textit{Let ${\mathbf C}_1,{\mathbf C}_2,{\mathbf C}_3$ be three abelian
categories, and suppose that in each of the first two every object is
isomorphic to a subobject of an injective object. Let $F$ and $G$ be covariant
functors from ${\mathbf C}_1$ to ${\mathbf C}_2$ and from ${\mathbf C}_2$ to ${\mathbf C}_3$, respectively, and
suppose that $G$ is left exact and that $F$ sends injective objects to $G$-acyclic objects. Finally, let
$\mathbf K$ be a covariant functor from ${\mathbf C}_1$ to the category
$\mathbf K({\mathbf C}_2)$ of nonnegative-degree complexes in ${\mathbf C}_2$, satisfying the two
conditions \textup{(i)} $H^0\mathbf K(A)=F(A)$ (a functorial isomorphism), and \textup{(ii)}
$\mathbf K(A)$ is acyclic in dimensions $n>0$ when $A$ is injective. Under
these conditions there exists on ${\mathbf C}_1$ a cohomological spectral functor
abutting to the graded functor $(R^n(GF)(A))$, whose initial term is}

\LayerTarget{tohoku:equation:5.5.1}
\[
\tag{5.5.1}
E_2^{p,q}(A)={\mathbf R}^pG\bigl(R^q\mathbf K(A)\bigr)
\]

\LayerTarget{tohoku:proof:lemma-5.5.1}
Let $H^0$ denote the covariant functor $L\to H^0(L)$ from
$K({\mathbf C}_2)$ to ${\mathbf C}_2$. We then have
$GF=G(H^0\mathbf K)=(GH^0)\mathbf K$. If $A$ is injective, $\mathbf K(A)$ is
$(GH^0)$-acyclic, i.e. ${\mathbf R}^nG(\mathbf K(A))=0$ for $n>0$; indeed,
$\mathbf K(A)$ is a resolution of $F(A)$ by (i) and (ii), so
${\mathbf R}^nG(\mathbf K(A))=R^nG(F(A))$, and the latter term is zero by the hypothesis
on $F$. We may therefore apply Theorem 2.4.1 to the composite functor
$(GH^0)\mathbf K$, which gives the spectral functor of the lemma.

\LayerTarget{tohoku:lemma:5.5.2}
\textsc{Lemma 5.5.2.} \textit{With ${\mathbf C}_1,{\mathbf C}_2,{\mathbf C}_3,F$, and $G$ as in
the preceding lemma, there exists a cohomological spectral functor on the
category $\mathbf K({\mathbf C}_1)$ of nonnegative-degree complexes $\mathrm C$ in ${\mathbf C}_1$, abutting
to the graded functor $({\mathbf R}^n(GF)(\mathrm C))$, whose initial term is}

\LayerTarget{tohoku:equation:5.5.2}
\[
\tag{5.5.2}
E_2^{p,q}(\mathrm C)={\mathbf R}^pG\bigl(R^qF(\mathrm C)\bigr).
\]

\LayerTarget{tohoku:proof:lemma-5.5.2:p1}
Place a bar over a functor to denote the functor obtained by ``extension'' to
the categories of nonnegative-degree complexes. We then have
${\mathbf R}^n(GF)=R^n\bigl(H^0\overline{GF}\bigr)=R^n\bigl((H^0\overline G)\overline F\bigr)$.
The functor $\overline F:\mathbf K({\mathbf C}_1)\to\mathbf K({\mathbf C}_2)$ sends
injective objects to $(H^0\overline G)$-acyclic objects; in other words, if
$\mathrm C$ is an injective object of $\mathbf K({\mathbf C}_1)$, then
${\mathbf R}^nG(\overline F(\mathrm C))=0$ for $n>0$; indeed, by the characterization

\TohokuContinuation{208}
